% BHU PYQ -- Manually transcribed & error-corrected
\documentclass[11pt,a4paper]{article}

\usepackage[a4paper,top=2.5cm,bottom=2.5cm,left=2.8cm,right=2.8cm,headheight=14pt]{geometry}
\usepackage{amsmath,amssymb}
\usepackage[T1]{fontenc}
\usepackage[utf8]{inputenc}
\usepackage{lmodern}
\usepackage{microtype}
\usepackage{fancyhdr}
\usepackage{enumitem}
\usepackage{hyperref}
\usepackage{lastpage}
\usepackage{parskip}

\hypersetup{colorlinks=true,urlcolor=black,linkcolor=black,
  pdftitle={BPE-601: Topics in Modern Physics -- BHU PYQ},pdfauthor={Banaras Hindu University}}

\pagestyle{fancy}\fancyhf{}
\renewcommand{\headrulewidth}{0.4pt}\renewcommand{\footrulewidth}{0.4pt}
\fancyhead[L]{\small   \textit{Banaras Hindu University}}
\fancyhead[R]{\small   \textit{BPE-601: Topics in Modern Physics}}
\fancyfoot[C]{\small Page  \thepage\ of \pageref{LastPage}}


\newlist{parts}{enumerate}{2}
\setlist[parts,1]{label=(\alph*),leftmargin=*,itemsep=5pt,topsep=3pt}
\setlist[parts,2]{label=(\roman*),leftmargin=*,itemsep=3pt,topsep=2pt}

\newcommand{\pts}[1]{\hfill   \textit{[#1]}}


\begin{document}


\begin{center}
  {\large\bfseries BANARAS HINDU UNIVERSITY}\\[2pt]
  {
\normalsize B.Sc. (Hons.) Semester VI Examination, 2022-23}\\[6pt]
  \rule{0.8\linewidth}{0.5pt}\\[6pt]
  {\large\bfseries Paper No. BPE-601: Topics in Modern Physics}\\[4pt]
  {
\normalsize Time: 3 Hours\hfill Maximum Marks: 70}
\end{center}
\rule{\linewidth}{0.4pt}
\smallskip



\noindent   \textit{Instructions: Answer all questions. Symbols have their usual meanings.}
\bigskip




\noindent   \textbf{Question 1.} Answer any   \textbf{four} of the following: \pts{4 $ \times$ 5 = 20}

\begin{parts}
  \item State the postulates of special theory of relativity, precisely. What assumption other than the relativity principle and constancy of $c$ has been used in deducing the Lorentz transformation equations?
  \item Write the Lorentz transformation matrix for Lorentz boost in the $x$-direction. Write down the tensor equations showing that this matrix is symmetric and its own inverse.
  \item A rocket ship is receding from Earth at a speed of $0.2c$. A light on the rocket ship of wavelength $\lambda_0 = 4500\, \text{\AA}$ appears blue to a passenger on the ship. What colour would it appear to an observer on Earth?
  \item Define astronomical unit, light year, and parsec.
  \item Consider two objects ($A$ and $B$) that are two astronomical units apart. The line joining $A$ and $B$ subtends an angle of $4''$ (4 arcsec) at a nearby star $S$. What is the distance of the star $S$ from the midpoint of the line $AB$?
  \item Explain why the luminosity of a Red Giant is more than the original star, despite having a lower effective temperature.
\end{parts}

\medskip\rule{\linewidth}{0.2pt}




\noindent   \textbf{Question 2.}

\begin{parts}
  \item Derive the relativistic equation for the aberration of light. Use this equation to show that the first-order aberration effect corresponds to the classical picture and explain the orientation of a telescope observing a star. \pts{7}
  \item Write down the conversion for electric charge from CGS to SI system of units. \pts{2}
  \item What is gauge freedom? Show that the electromagnetic tensor $F_{\mu
u}$ is invariant under the gauge transformation $A_\mu  o A_\mu + \partial_\mu \chi$, where $\chi$ is a scalar. \pts{3.5}
\end{parts}

\medskip\rule{\linewidth}{0.2pt}




\noindent   \textbf{Question 3.}

\begin{parts}
  \item Derive the transformation equations for the electric field under special relativity and use them to deduce the transformation for the magnetic field using the correspondence between the two. \pts{6.5}
  \item The average lifetime of a $\pi$ meson in its own frame of reference is $26\, \text{ns}$. If the $\pi$ meson moves with speed $0.95c$ with respect to the Earth, what is its lifetime as measured by an observer at rest on Earth? What is the average distance it travels before decaying as measured by an observer at rest on Earth? \pts{6}
\end{parts}
\hfill   \textbf{OR}

\begin{parts}
  \item Use the concept of light cone to explain that violating the speed of light is a must to violate causality. \pts{7}
  \item Write Maxwell's equations in the covariant form using the electromagnetic tensor and its Hodge dual. Label the equations and explain the terms on the RHS. \pts{5.5}
\end{parts}

\medskip\rule{\linewidth}{0.2pt}




\noindent   \textbf{Question 4.}

\begin{parts}
  \item Derive the relation between the thermal and gravitational energy considering hydrostatic equilibrium for a star. What are the consequences of this relationship? \pts{12.5}
\end{parts}
\hfill   \textbf{OR}

\begin{parts}
  \item Explain penetration probability in reference to the nuclear fusion reaction at the core of a star. \pts{5}
  \item Compare the different energy generation mechanisms taking place at the core of stars. Discuss the main-sequence lifetime of stars. \pts{7.5}
\end{parts}

\medskip\rule{\linewidth}{0.2pt}




\noindent   \textbf{Question 5.}

\begin{parts}
  \item Discuss at least two methods by which dark matter may be observationally detected. \pts{8}
  \item Explain the working of LIGO used for the detection of gravitational waves. \pts{4.5}
\end{parts}

\vfill
\rule{\linewidth}{0.4pt}

\begin{center}\small
  \href{https://bhu-exam.vercel.app/}{Click here to visit our website}\\[4pt]
  \href{https://me-aryan.vercel.app/}{Click here to visit our contributor}\\[4pt]
  \href{https://quashan-me.vercel.app/}{Click here to visit our contributor}
\end{center}

\end{document}
